%% kapa: exam=1 round=33 year=2022 period=2 form=A subject=appraisal_law subject_ko=감정평가관계법규 seq=24 number=24
%% source: Q-Net 감정평가사 기출문제 (cst00309) artlSeq=5209997 file=raw/1008/5209997/2022년도 제33회 감정평가사 1차시험 2교시(감정평가 관계법규, 회계학).pdf
%% conversion: kordoc 4.24.51 -> tools/gen1.py (auto); review=proofread-sonnet; flags=none
\begin{question}{24}
건축법상 용어의 정의에 관한 조문의 일부이다. ( )에 들어갈 내용으로 옳은 것은?
\begin{framedtext}
○ “지하층”이란 건축물의 바닥이 지표면 아래에 있는 층으로서 바닥에서 지표면까지 평균높이가 해당 층 높이의 ( ㄱ )분의 1 이상인 것을 말한다. \par
○ “고층건축물”이란 층수가 ( ㄴ )층 이상이거나 높이가 ( ㄷ )미터 이상인 건축물을 말한다. \par
\end{framedtext}
\begin{choices}
  \item ㄱ: 2, ㄴ: 20, ㄷ: 100
  \item ㄱ: 2, ㄴ: 20, ㄷ: 120
  \item ㄱ: 2, ㄴ: 30, ㄷ: 120
  \item ㄱ: 3, ㄴ: 20, ㄷ: 100
  \item ㄱ: 3, ㄴ: 30, ㄷ: 120
\end{choices}
\end{question}
\answer{3}
